\section{Compact kernels, Fredholm determinants, and finite model spaces}
\label{sec:hilbert-schmidt-fredholm}
% BEGIN CURATED SUBJECT INDEX
\index{spectrum!point}
\index{resolvent}
\index{direct integral}
\index{compact operator}
\index{Fredholm theory}
\index{Fredholm determinant}
\index{Lippmann--Schwinger equation}
\index{resonance!residue}
\index{pseudostate}
\index{CDCC}
% END CURATED SUBJECT INDEX
\index{compact operator}\index{Fredholm determinant}\index{Galerkin method}

\calculationroute
  {The operator topologies, direct integrals, and resolvents constructed in
  Section~\ref{sec:functional-operator-algebra}.}
  {Identify the compact operator hidden in a scattering equation, replace
  the infinite-dimensional inverse by a Fredholm determinant, and calculate
  the residue of a simple pole.}
  {Know exactly why a finite-basis calculation can converge even though the
  identity operator itself cannot be approximated in operator norm by
  finite-rank projections.}

The Lippmann--Schwinger equation looks like a linear equation,
\begin{equation}
  (\identity-K(z))\psi(z)=\phi(z),
  \label{eq:fredholm-start}
\end{equation}
but an inverse of an infinite matrix need not share the elementary properties
of a finite inverse.  The missing hypothesis is usually compactness.  Once
the interaction has been placed in a compact kernel, three operations used
throughout this book become legitimate versions of one another:
\begin{equation}
  \begin{aligned}
    \text{solve an integral equation}
    &\quad\Longleftrightarrow\quad
    \text{locate a determinant zero},\\
    &\quad\Longleftrightarrow\quad
    \text{extract a finite-rank pole residue}.
  \end{aligned}
  \label{eq:fredholm-three-operations}
\end{equation}
This section develops that statement from the definition of a compact
operator.  The theory is also the mathematical reason why Gaussian
pseudostates, Galerkin bases, and CDCC model spaces can approximate tested
continuum amplitudes without turning every positive-energy eigenvalue into a
physical level.  Standard references include
Refs.~\cite{Kato:1995,SimonTraceIdeals:2005,Kukulin:1989}.

\subsection{Compactness is uniform finite-dimensionality after the operator acts}
\label{subsec:compact-operators}
% BEGIN CURATED SUBJECT INDEX
\index{Hilbert space}
\index{spectrum!point}
\index{Riesz projection}
\index{compact operator}
\index{Jordan chain}
% END CURATED SUBJECT INDEX

A bounded operator $K:\Hil_1\to\Hil_2$ is compact if it maps the unit ball of
$\Hil_1$ into a relatively compact subset of $\Hil_2$.  Equivalently, for
every bounded sequence $\{\psi_n\}$ there is a subsequence for which
$K\psi_{n_j}$ converges in norm.  A finite-rank operator is compact, and a
norm limit of finite-rank operators is compact.  On Hilbert space the
converse also holds: every compact operator is a norm limit of finite-rank
operators.

This definition contains the important order of operations.  Let $P_N$ be
orthogonal projections onto an increasing sequence of finite-dimensional
spaces with $P_N\to\identity$ strongly.  In an infinite-dimensional space,
\begin{equation}
  \|\identity-P_N\|=1
  \label{eq:projection-not-norm-identity}
\end{equation}
for every finite $N$.  Thus a basis truncation never approximates the
identity in operator norm.  If $K$ is compact, however,
\begin{equation}
  \|(\identity-P_N)K\|\longrightarrow0,
  \qquad
  \|K(\identity-P_N)\|\longrightarrow0.
  \label{eq:compact-projection-norm}
\end{equation}
To prove the first limit, cover the compact closure of the image of the unit
ball by finitely many small balls.  Strong convergence is uniform on their
finite set of centers and hence, by the triangle inequality, on the whole
compact image.  The second statement follows by applying the first to
$K^\dagger$.  Consequently
\begin{equation}
  K_N=P_NKP_N
  \quad\Longrightarrow\quad
  \|K_N-K\|\to0 .
  \label{eq:galerkin-compact-convergence}
\end{equation}
Equation~\eqref{eq:galerkin-compact-convergence}, not an imagined norm
convergence $P_N\to\identity$, is the compactness mechanism behind a stable
Galerkin calculation.

The nonzero spectrum of a compact operator consists of isolated eigenvalues
of finite algebraic multiplicity, with zero as the only possible accumulation
point.  A nonzero spectral point therefore has a finite-rank Riesz projection.
For a non-normal $K$, right and left eigenvectors need not be orthogonal and
Jordan chains may occur, but the root space belonging to a nonzero isolated
eigenvalue remains finite dimensional.  This is the first place where the
infinite-dimensional problem recovers finite matrix algebra without being
replaced by it.

\subsection{Hilbert--Schmidt and trace-class kernels}
\label{subsec:schatten-ideals}
% BEGIN CURATED SUBJECT INDEX
\index{spectrum!point}
\index{resolvent}
\index{compact operator}
\index{Hilbert--Schmidt operator}
\index{trace-class operator}
% END CURATED SUBJECT INDEX
\index{Hilbert--Schmidt operator}\index{trace-class operator}
\index{Schatten ideal}

Let $s_n(K)$ be the singular values of a compact operator, namely the
eigenvalues of $|K|=(K^\dagger K)^{1/2}$ in decreasing order and with
multiplicity.  The Schatten ideal $\mathfrak S_p$ consists of operators for
which
\begin{equation}
  \|K\|_p^p=\sum_n s_n(K)^p<\infty .
  \label{eq:schatten-norm}
\end{equation}
The two classes needed here are
\begin{align}
  \mathfrak S_2&=\{\text{Hilbert--Schmidt operators}\},
  &\|K\|_2^2&=\tr(K^\dagger K),
  \label{eq:hilbert-schmidt-definition}\\
  \mathfrak S_1&=\{\text{trace-class operators}\},
  &\|K\|_1&=\tr|K|.
  \label{eq:trace-class-schatten-definition}
\end{align}
One has $\mathfrak S_1\subset\mathfrak S_2\subset\mathcal K(\Hil)$, where
$\mathcal K(\Hil)$ denotes the compact operators.  If $A$, $B$ are bounded and
$K\in\mathfrak S_p$, then
\begin{equation}
  AKB\in\mathfrak S_p,
  \qquad
  \|AKB\|_p\leq\|A\|\,\|K\|_p\,\|B\|.
  \label{eq:schatten-ideal-property}
\end{equation}
Thus these are two-sided ideals in $\BHil$.

For $\Hil=L^2(X,\rmd\mu)$, an integral kernel $K(x,y)$ satisfying
\begin{equation}
  \int_X\int_X |K(x,y)|^2\dd\mu(x)\rmd\mu(y)<\infty
  \label{eq:hs-kernel-criterion}
\end{equation}
defines a Hilbert--Schmidt operator, and the integral is $\|K\|_2^2$.
This elementary criterion is the most useful entry point for scattering.
It also explains why the order in which potential factors and a resolvent are
multiplied matters: a free resolvent is not compact on all of space, but a
resolvent localized on both sides by a sufficiently short-range potential can
be Hilbert--Schmidt.

\begin{warningbox}
Neither a delta function nor a continuous-spectrum identity kernel is square
integrable on $X\times X$.  Discretizing such a kernel produces a finite
matrix, but finiteness of that matrix is not evidence that the original
operator was Hilbert--Schmidt.  Compactness must be checked before the
discretization or proved in a weighted realization.
\end{warningbox}

\subsection{The Birman--Schwinger kernel of a Schr\"odinger operator}
% BEGIN CURATED SUBJECT INDEX
\index{spectrum!continuous}
\index{spectrum!point}
\index{resolvent}
\index{compact operator}
\index{Fredholm theory}
\index{Birman--Schwinger operator}
\index{Jost function}
\index{boundary condition!incoming and outgoing}
\index{scattering!Coulomb}
% END CURATED SUBJECT INDEX
\index{Birman--Schwinger operator}
\label{subsec:birman-schwinger}

Let $H=H_0+V$ on $L^2(\real^d)$ and use the resolvent convention
\begin{equation}
  \Resolvent_0(z)=(z-H_0)^{-1}.
  \label{eq:bs-free-resolvent}
\end{equation}
For a real multiplication potential write
$V=\sgn(V)|V|$ and define
\begin{equation}
  K(z)=|V|^{1/2}\Resolvent_0(z)\sgn(V)|V|^{1/2}.
  \label{eq:birman-schwinger-kernel}
\end{equation}
If $(z-H)\psi=0$ and $\phi=|V|^{1/2}\psi$, then
\begin{equation}
  (\identity-K(z))\phi=0.
  \label{eq:birman-schwinger-equation}
\end{equation}
Conversely, a solution of Eq.~\eqref{eq:birman-schwinger-equation} gives
$\psi=\Resolvent_0(z)\sgn(V)|V|^{1/2}\phi$, subject to the usual
domain qualifications.  Hence an eigenvalue or continued pole of $H$ occurs
when $1$ is an eigenvalue of $K(z)$.  If one instead defines the resolvent as
$(H_0-z)^{-1}$, the corresponding eigenvalue is $-1$; many sign differences
in the literature are only this convention.

For $d=3$ and $z$ away from $[0,\infty)$, the free kernel is proportional to
\begin{equation}
  G_0(z;\bm x-\bm y)
  =\frac{\rme^{-\kappa(z)|\bm x-\bm y|}}
  {|\bm x-\bm y|},
  \qquad \re\kappa(z)>0,
  \label{eq:off-shell-free-green-decay}
\end{equation}
up to constants fixed in Section~\ref{subsec:free-resolvent-ls}.  If, for
example, $V$ is bounded and compactly supported, then
\begin{equation}
  \|K(z)\|_2^2
  \propto
  \iint
  |V(\bm x)|,|G_0(z;\bm x-\bm y)|^2,|V(\bm y)|
  \dd^3x\rmd^3y<\infty.
  \label{eq:bs-hs-integral}
\end{equation}
The singularity $|\bm x-\bm y|^{-2}$ is locally integrable in three
dimensions, while the localization controls large separations.  Much weaker
hypotheses are available, but this sufficient condition displays the
mechanism without hiding it in a theorem.

On the continuous spectrum the exponential decay in
Eq.~\eqref{eq:off-shell-free-green-decay} disappears.  Compactness and
boundary values are then formulated between weighted spaces, or after
localization by the potential.  A Coulomb tail changes the comparison
resolvent itself.  Thus the Birman--Schwinger argument does not license using
the free kernel for the long-range problems of
Section~\ref{sec:long-range-ir}.

The resolvent identity follows by solving Eq.~\eqref{eq:fredholm-start}:
\begin{align}
  \Resolvent(z)
  &=\Resolvent_0(z)
  +\Resolvent_0(z)\sgn(V)|V|^{1/2}
  [\identity-K(z)]^{-1}
  |V|^{1/2}\Resolvent_0(z).
  \label{eq:bs-resolvent-formula}
\end{align}
Every nontrivial singularity of the right-hand side is therefore carried by
$[\identity-K(z)]^{-1}$.  This is the operator whose determinant will be
compared with a Jost function and with the exact-WKB incoming coefficient.

\subsection{Analytic Fredholm theorem: why the continued inverse is meromorphic}
% BEGIN CURATED SUBJECT INDEX
\index{spectrum!point}
\index{resolvent}
\index{compact operator}
\index{Fredholm theory}
% END CURATED SUBJECT INDEX
\index{analytic Fredholm theorem}\index{meromorphic resolvent}
\label{subsec:analytic-fredholm}

Let $\Omega\subset\complex$ be connected and let
$z\mapsto K(z)$ be analytic in operator norm with compact values.  The
analytic Fredholm theorem states that exactly one of the following occurs:
\begin{enumerate}[label=(\roman*)]
  \item $\identity-K(z)$ is noninvertible for every $z\in\Omega$;
  \item $[\identity-K(z)]^{-1}$ exists except at a discrete set and is a
  meromorphic operator-valued function whose principal parts have finite
  rank.
\end{enumerate}
If the inverse exists at one point, alternative (ii) holds.  In scattering,
large negative energy or sufficiently small coupling usually supplies that
one point.

The finite-dimensional content of the theorem can be seen locally.  Near a
point $z_0$ where $1$ is an isolated eigenvalue of $K(z_0)$, take a small
spectral contour around that eigenvalue and form its finite-rank Riesz
projection $P$.  Decompose $\Hil=P\Hil\oplus(\identity-P)\Hil$.  The block on
$(\identity-P)\Hil$ is invertible near $z_0$; a Schur complement reduces the
remaining obstruction to a finite matrix on $P\Hil$.  Its ordinary
determinant has isolated zeros, and its inverse has a finite-rank Laurent
principal part.  Compactness is precisely what made the obstructing block
finite.

This proof also identifies a limitation.  A threshold is not, in general, an
ordinary interior point of an analytic domain: it is a branch point of the
continued kernel.  The theorem is applied on a chosen momentum sheet or in a
uniformizing variable, not blindly across the threshold.

\subsection{Fredholm and regularized determinants}
\label{subsec:fredholm-determinants}
% BEGIN CURATED SUBJECT INDEX
\index{spectrum!point}
\index{Hilbert--Schmidt operator}
\index{trace-class operator}
\index{Fredholm theory}
\index{Fredholm determinant}
\index{Jost function}
\index{connection coefficient}
\index{boundary condition!incoming and outgoing}
% END CURATED SUBJECT INDEX

If $K(z)$ is trace class, its Fredholm determinant is
\begin{equation}
  D(z)=\det(\identity-K(z))
  =\prod_n[1-\lambda_n(z)],
  \label{eq:fredholm-determinant}
\end{equation}
where the eigenvalues are counted with algebraic multiplicity.  The product
converges, $D$ is analytic, and
\begin{equation}
  D(z)=0
  \quad\Longleftrightarrow\quad
  \identity-K(z)\text{ is not invertible}.
  \label{eq:fredholm-zero}
\end{equation}
Where the inverse exists,
\begin{equation}
  \frac{\rmd}{\rmd z}\log D(z)
  =-\tr\!\left([\identity-K(z)]^{-1}K'(z)\right).
  \label{eq:fredholm-log-derivative}
\end{equation}

A Hilbert--Schmidt operator need not be trace class, so the linear term in
the product can diverge.  The regularized determinant
\begin{equation}
  \det{}_2(\identity-K)
  =\det\bigl[(\identity-K)\rme^K\bigr]
  =\prod_n(1-\lambda_n)\rme^{\lambda_n}
  \label{eq:det2-definition}
\end{equation}
is well defined for $K\in\mathfrak S_2$.  The exponential removes the
non-summable linear contribution.  Most importantly,
\begin{equation}
  \det{}_2(\identity-K)=0
  \quad\Longleftrightarrow\quad
  1\in\spectrum(K),
  \label{eq:det2-zero}
\end{equation}
with the same zero multiplicity.  Multiplication of a resonance determinant
by a nonvanishing analytic factor changes neither its zeros nor the local
null spaces.  This is why a Jost function, a Fredholm determinant, and an
exact-WKB connection coefficient may differ in normalization while encoding
the same poles.

For a partial-wave Schr\"odinger problem one may schematically write
\begin{equation}
  \Cincoming(E)
  =N(E)F_l(k(E))
  =\widetilde N(E)\det{}_2[\identity-K_l(E)],
  \label{eq:cin-jost-fredholm}
\end{equation}
on a domain where all three objects have been constructed and where
$N$, $\widetilde N$ do not vanish.  Equation~\eqref{eq:cin-jost-fredholm} is not
a universal equality of preferred normalizations.  It is the precise local
statement that the three analytic coordinates define the same divisor of
resonance zeros.

\subsection{Simple-zero residue without diagonalizing the whole operator}
% BEGIN CURATED SUBJECT INDEX
\index{resolvent}
\index{Fredholm theory}
\index{scattering!multichannel}
\index{connection coefficient}
\index{boundary condition!incoming and outgoing}
\index{resonance!residue}
\index{exact WKB}
\index{Feshbach projection}
% END CURATED SUBJECT INDEX
\index{residue}\index{left eigenvector}\index{right eigenvector}
\label{subsec:fredholm-simple-residue}

Set $A(z)=\identity-K(z)$ and suppose $\vsub{z}{R}$ is a simple
noninvertibility point.  Choose right and left null vectors
\begin{equation}
  A(\vsub{z}{R})u=0,
  \qquad
  A(\vsub{z}{R})^\dagger v=0,
  \label{eq:fredholm-left-right-null}
\end{equation}
where the second equation represents the left functional in Hilbert-space
notation.  If
\begin{equation}
  \vsub{\mathcal N}{R}
  =\dualpair{v}{A'(\vsub{z}{R})u}\ne0,
  \label{eq:fredholm-derivative-normalization}
\end{equation}
then
\begin{equation}
  A(z)^{-1}
  =\frac{\ket{u}\bra{v}}
  {(z-\vsub{z}{R})\vsub{\mathcal N}{R}}
  +\vsub{A}{reg}(z).
  \label{eq:fredholm-inverse-residue}
\end{equation}
The result follows by applying $A(z)=A(\vsub{z}{R})+
(z-\vsub{z}{R})A'(\vsub{z}{R})+\dots$ to the proposed Laurent term.
The scale changes $u\mapsto cu$, $v\mapsto c^{-1*}v$ cancel between numerator
and denominator.

Substitution in Eq.~\eqref{eq:bs-resolvent-formula} gives the factorized
resolvent residue.  The denominator
$\dualpair{v}{A'u}$ is the operator form of every derivative
normalization encountered later:
\begin{equation}
  \dualpair{v}{A'u}
  \quad\longleftrightarrow\quad
  F_l'(k_{\mathrm R})
  \quad\longleftrightarrow\quad
  \Cincoming'(\resonantE)
  \quad\longleftrightarrow\quad
  1-\Sigma'(\resonantE).
  \label{eq:derivative-normalization-dictionary}
\end{equation}
The last expression occurs after Feshbach elimination.  Neglecting the energy
derivative of an effective interaction is therefore not an innocent
normalization choice; it changes the pole residue.

For a matrix incoming block $\Cincoming(E)$, the identical
calculation gives
\begin{equation}
  \Cincoming(E)^{-1}
  \sim
  \frac{\ket{u}\bra{v}}
  {(E-\resonantE)
   \dualpair{v}{\partial_E\Cincoming(\resonantE)u}}.
  \label{eq:matrix-connection-residue}
\end{equation}
This is the multichannel extension of dividing by
$\Cincoming'(\resonantE)$ in exact WKB.

\subsection{Worked model: a rank-one interaction}
\label{subsec:rank-one-fredholm-model}
% BEGIN CURATED SUBJECT INDEX
\index{spectrum!point}
\index{resolvent}
\index{Fredholm theory}
\index{connection coefficient}
\index{boundary condition!incoming and outgoing}
\index{resonance!residue}
% END CURATED SUBJECT INDEX

Let
\begin{equation}
  V=\lambda\ket{g}\bra{g},
  \qquad g\in\Hil .
  \label{eq:rank-one-potential}
\end{equation}
The transition operator has the separable form
\begin{equation}
  T(z)=\ket{g}\tau(z)\bra{g}.
  \label{eq:rank-one-T-ansatz}
\end{equation}
Inserting it in $T=V+V\Resolvent_0T$ yields
\begin{align}
  \tau(z)&=\frac{\lambda}{D(z)},
  &D(z)&=1-\lambda g_0(z),
  &g_0(z)&=\braket{g|\Resolvent_0(z)|g}.
  \label{eq:rank-one-denominator}
\end{align}
The entire pole problem is now one Fredholm eigenvalue.  If
$D(\vsub{z}{R})=0$ and $D'(\vsub{z}{R})\ne0$, then
\begin{equation}
  \Res_{z=\vsub{z}{R}}T(z)
  =\frac{\lambda\ket{g}\bra{g}}{D'(\vsub{z}{R})}
  =-\frac{\ket{g}\bra{g}}{g_0'(\vsub{z}{R})}.
  \label{eq:rank-one-residue}
\end{equation}
The first equality has exactly the form
$\mathcal N(\resonantE)/\Cincoming'(\resonantE)$.  A numerical
calculation that reproduces $\vsub{z}{R}$ but not $g_0'(\vsub{z}{R})$
has not yet reproduced the resonance strength.

The discontinuity of $g_0(E\pm\rmi0)$ supplies the open-channel width.  A
threshold branch point remains in $g_0(z)$ even though the algebra in
Eq.~\eqref{eq:rank-one-denominator} is one dimensional.  Finite rank does not
remove continuum analyticity; it isolates where that analyticity enters.

\subsection{Galerkin, pseudostates, and CDCC as compact-kernel quadrature}
% BEGIN CURATED SUBJECT INDEX
\index{spectrum!point}
\index{resolvent}
\index{Riesz projection}
\index{Hilbert--Schmidt operator}
\index{Fredholm theory}
\index{Fredholm determinant}
\index{resonance!residue}
\index{pseudostate}
\index{CDCC}
% END CURATED SUBJECT INDEX
\index{Galerkin method}\index{pseudostate}\index{CDCC}
\label{subsec:fredholm-galerkin-cdcc}

Let $K_N=P_NKP_N$.  If $K$ is compact, Eq.~\eqref{eq:galerkin-compact-convergence}
implies that isolated nonzero eigenvalues and their total Riesz projections
are stable under a sufficiently large truncation.  If $K\in\mathfrak S_2$
and $\|K_N-K\|_2\to0$, the regularized determinants converge locally
uniformly.  A practical pole calculation may therefore proceed as follows:
\begin{enumerate}[label=\arabic*.]
  \item construct the compact or Hilbert--Schmidt kernel before projection;
  \item choose $P_N$ and calculate the finite matrix $P_NK(z)P_N$;
  \item locate zeros of its determinant on a declared sheet;
  \item enlarge $P_N$ and track a whole Riesz cluster, not a single sorted
  eigenvalue;
  \item calculate the left--right derivative denominator in
  Eq.~\eqref{eq:fredholm-derivative-normalization};
  \item test the resulting residue between the physical source and detector.
\end{enumerate}

The same logic clarifies continuum pseudostates.  The projection $P_N$ is a
finite quadrature for matrix elements of a compactly tested resolvent.  Most
positive-energy eigenvalues depend on basis ranges and should move as the
quadrature is refined.  A resonance is identified by stability of a
continued pole or of its Riesz projection and residue.  CDCC adds coupled
relative motion and physical channel boundary conditions, but its internal
pseudostate space obeys the same rule.

\subsection{Compact kernels versus decomposable scattering operators}
% BEGIN CURATED SUBJECT INDEX
\index{measure theory}
\index{resolvent}
\index{direct integral}
\index{decomposable operator}
\index{compact operator}
\index{Hilbert--Schmidt operator}
\index{Fredholm theory}
\index{Birman--Schwinger operator}
\index{scattering matrix}
\index{scattering!on-shell}
% END CURATED SUBJECT INDEX
\index{direct integral}\index{decomposable operator}
\label{subsec:compact-versus-direct-integral}

Compactness should not be attached to the wrong operator.  In the spectral
representation of $H_0$, the scattering operator is decomposable,
\begin{equation}
  \Scattering=\int^\oplus \Scattering(E)\dd E.
  \label{eq:S-decomposable-repeat}
\end{equation}
On a non-atomic energy interval, a nonzero decomposable multiplication
operator is ordinarily not compact.  Indeed, if
$\|\Scattering(E)-\identity\|>\varepsilon$ on a set of positive measure,
divide that set into disjoint subsets and choose unit sections supported on
them.  Their images remain separated, so no convergent subsequence exists.

There is no contradiction.  The compact object is the localized off-shell
kernel, a Birman--Schwinger operator, or a difference of suitably weighted
resolvents.  The on-shell scattering operator is a measurable family acting
throughout a continuum.  Direct-integral theory tells us how it is organized;
Fredholm theory tells us how its analytic singularities arise.  Exact WKB
then calculates the same singularity from global connection data.

\begin{checkpointbox}
Before proceeding, the reader should be able to distinguish the strong limit
$P_N\to\identity$ from the norm limit $P_NKP_N\to K$, verify the
Hilbert--Schmidt criterion for a localized Green kernel, derive
Eq.~\eqref{eq:fredholm-inverse-residue}, and explain why the derivative of a
determinant or incoming coefficient is part of the physical pole residue.
\end{checkpointbox}
